\section{Keep products and partial sums exact}
The storage width of an input is four bytes, while its permitted numerical range
is signed sixteen-bit. Load it with \texttt{lw}, which sign-extends a word into
the full register. Thus stored \(-32768\) becomes the 64-bit signed integer
\(-32768\). An unsigned word load would instead turn its stored bit pattern into
4294934528, a different input value.

The largest positive product is
\[
  (-32768)(-32768)=2^{30}=1073741824.
\]
One product fits in a signed word. A sum of sixteen such products is
\(2^{34}=17179869184\), which needs thirty-six signed bits. The negative endpoint
is \(-17179344896\). Every partial sum of at most sixteen permitted products
also fits in a signed doubleword. Use a 64-bit accumulator throughout the reduction
and store all eight output bytes.

The illustrated implementation uses \texttt{dmul} for the product and
\texttt{daddu} for addition. The lack of a signed-overflow trap in \texttt{daddu}
does not imply a loss of accuracy: the proven domain keeps all relevant sums inside
the signed 64-bit range. A native signed-word \texttt{mul} whose result is sign-extended
also produces the exact individual product under this input bound. Changing the
product instruction can therefore be a legal equivalent. Narrowing the accumulator
after each addition is a different change and fails the two-product boundary.

\begin{table}[htbp]
  \centering
  \caption{A boundary reduction with both inputs equal to \(-32768\).}
  \tagpdfsetup{table/header-rows={1}}
  \begin{tabular}{rrll}
    \toprule
    Products & Exact sum   & High word         & Low word          \\
    \midrule
    1        & 1073741824  & \texttt{00000000} & \texttt{40000000} \\
    2        & 2147483648  & \texttt{00000000} & \texttt{80000000} \\
    4        & 4294967296  & \texttt{00000001} & \texttt{00000000} \\
    16       & 17179869184 & \texttt{00000004} & \texttt{00000000} \\
    \bottomrule
  \end{tabular}
\end{table}

\begin{MARSExercise}{Find the first lost distinction}
  A faulty implementation uses a signed word accumulator, then sign-extends its
  final result into the output doubleword. Trace its second addition for this table.
  What high word does it store, and what high word is required? Explain why widening
  only at the final store cannot repair the sum.
\end{MARSExercise}
